\scanpage{186}
Finalement, on trouve

\uline{Proposition} Soit $S$ un anneau commutatif,
$\mathcal A_{\unsure{1}}=S\{X_0,dX_0\}$, \deleted{\unsure{on cherche à}} $S$-algèbre
diff. à puiss. divisées, avec augmentation
\[
\alpha:\mathcal A_1\mathbin{\unsure{-|\longrightarrow}}S\qquad
(\alpha(X_0^{(i)})=0\quad\forall i\geq0,\ \alpha(dX_0)=0).
\]
[\uline{NB} on regarde \unsure{bien sûr} $S$ comme anneau
différentiel à p.d. où $d_S=0$, idéal nul pour
les p.d., et $\alpha$ est compatible avec ces structures)

On cherche les endomorphismes \inserted{$\varphi$} d'\inserted{$S$-}\uline{algèbre}
\uline{différentielle} \uline{augmentée} (mais sans
p.d. !) $\varphi$ de \deleted{$\unsure{S}$} $\mathcal A$
\[
\varphi:\mathcal A\longrightarrow\mathcal A\qquad\text{endom. de $S$-algèbre}
\]
\margin{$\deleted{\varphi d\unsure{=d}}$ \inserted{\unsure{puis.}}\quad
$\varphi(x+y)=\unsure{\varphi x)|\varphi(x)}+\varphi(y)$\par
$\varphi(xy)=\varphi(x)\varphi y$\par
$\unsure{\varphi(1)|\varphi1}=1$}
\[
\varphi d=d\varphi
\qquad\qquad
\alpha\varphi=\alpha
\]
Alors $\varphi$ est donné par les
\[
\begin{array}{ccl}
\varphi(X_0^{(i)})=\xi_i\in\mathcal A&&(i\geq1)\\
\text{\textquotedbl}&&\\
A_i+B_i\,dX_0&&A_i,B_i\in S\{X_0\}\ (i\geq1)
\end{array}
\]
(on aura $\varphi(dX_0)=\deleted{\unsure{\varphi dX}}d\varphi(X_0)=d\xi_1=dA_1=\deleted{\unsure{\partial}}A'_1dX_0$)
avec les conditions
\[
\begin{cases}
A_i=A_1^{(i)}\\
\alpha(A_1)=0\quad\text{i.e.}\quad A_1\in S\{X_0\}^{+}
\end{cases}
\]
donc $\varphi$ est donné par
\[
\boxed{\boxed{\begin{cases}
A_1\in S\{X_0\}^{+}\\
B_i\in S\{X_0\}\quad(i\geq1)
\end{cases}}}
\]
sans autre condition !

Pour que $\varphi$ commute aux p.d. il f. et s.
que l'on ait
\[
\xi_i=\xi_1^{(i)}\qquad(i\geq2)
\]
ou encore
\[
B_i=A_1^{(i-1)}B_1
\]
\margin{185}
